\documentclass[../mainReport.tex]{subfiles}
\begin{document}
\section{Report2}
\subsection{回文数组}
\subsubsection{形式化定义}
将一个数字序列经过若干次操作后变成回文序列, 求出将该数组变为回文数组所需的最小操作次数。操作支持如下：
\begin{itemize}
    \item 选择数组中相邻的两个数，同时将其加1或减1；
    \item 选择数组中单独的一个数，将其加1或减1。
\end{itemize}

\subsubsection{解题思路}
联想这道题的解题思路：\url{https://www.luogu.com.cn/problem/P5019}

我们可以想到\textcolor{blue}{贪心}的思路，上面那道题中，我们可以靠前一个坑\textcolor{blue}{顺便}填充下一个坑，这道题我们可以发现第1个操作就是一个\textcolor{blue}{顺便}的操作。所以我们记录好每个位置的\textcolor{blue}{对称差}，再让前一个对称差能够顺便填充后一个对称差即可。\textcolor{blue}{不过需要注意的是对称差是有正负的，和坑洞填充只有一个方向是不同的}。

% \subsubsection{完整代码}
% \begin{lstlisting}
% #include <iostream>
% #include <cstdio>
% #include <cmath>
% #define ll long long
% using namespace std;

% const int N = 2e5 + 10;

% int main() {
%     int n;
%     scanf("%d", &n);
%     int a[N];
%     for (int i = 1; i <= n; ++i) {
%         scanf("%d", a + i);
%     }

%     ll tot = 0;
%     for (int i = 1; i <= n / 2; ++i) {
%         int diff = a[i] - a[n - i + 1];
%         tot += abs(diff);

%         if (i + 1 <= n / 2) {
%             int next_diff = a[i + 1] - a[n - i];
%             if ((diff >= 0 && next_diff >= 0) || (diff <= 0 && next_diff <= 0)) {
%                 if (abs(diff) >= abs(next_diff)) {
%                     a[i + 1] = a[n - i];
%                 } else {
%                     a[i + 1] -= diff;
%                 }
%             }
%         }
%     }
%     printf("%lld", tot);
%     return 0;
% }
% \end{lstlisting}

\subsection{平衡饮食}
\subsubsection{形式化定义}
背包的容量为$X$，其中要装入3种物品，每种物品有其对应的价值$A_i$和体积$C_i$，现在希望你采取合适的装包方式，使得3种物品各自的总价值$F_1,F_2,F_3$的最小值最大。即，使得$\min(F_1,F_2,F_3)$最大。

\subsubsection{解题思路}
采用\textbf{动态规划}和\textbf{二分答案}的思想。

\begin{itemize}
    \item 状态定义：$dp[i][j]$，种类$i$用$j$容量所能获得的价值最大值
    \item 状态转移方程：$dp[i][j]=\max(dp[i][j], dp[i][j-c] + a)$，对每种食物，需要逆向更新数组，经典的01背包问题。
    \item 初始条件：$dp[i][0] = 0$，其他均设置为$INF$，表示该情况非法。 
    \item 前缀最大值优化：$dp[i][j] = \max(dp[i][j], dp[i-1][j])$，如果当前容量$j$下，无法获得价值，则直接取上一步的最大值。
\end{itemize}

对于本题，显然是满足单调性的，如果可以获取更大的最低价值，则一定可以获取更小的最低价值。所以满足\textcolor{blue}{二分答案}的条件。即，将优化问题转化为判定问题。对于给定的目标值$mid$，判断是否存在一个组合可以使得在满足容量限制的情况下，3种物品的总价值均大于$mid$。这样就能说明答案至少是$mid$，可以尝试更大的的答案。

% \begin{lstlisting}[caption={二分答案部分代码}]
% // 二分查找
% int low=0, high=min({dp[1][x], dp[2][x], dp[3][x]});
% while(low <= high) {
%     int mid = (low+high)/2;
%     if(可以满足min(V1,V2,V3)>=mid) 
%         low = mid+1;
%     else 
%         high = mid-1;
% }

% // 二分验证mid
% int sum = 0, ok = 1;
% for (int v = 1; v <= 3; v++) {
%     int cost = lower_bound(dp[v], dp[v] + x + 1, mid) - dp[v];
%     if ((sum += cost) > x || cost > x) {
%         ok = 0; break;
%     }
% }
% \end{lstlisting}

\subsubsection{复杂度分析}
对于状态转移部分，每个物品都要遍历$x-c$次，总共$n$个物品，所以状态函数计算部分的复杂度为：
\[
    T_{DP}=O(n\cdot x)
\]
二分答案对于验证部分需要对3种物品进行二分查找：
\[
    T_{check}=O(3\cdot\log x)=O(\log x)
\]

设物品中最大的价值为$V_{max}$，则二分的次数为：
\[
    T_{search}=O(\log V_{max})
\]

所以总时间复杂度为：
\[
    T_{total}=T_{DP} + T_{check}\cdot T_{search}=O(n\cdot x) + O(\log x\cdot\log V_{max})
\]

% \subsubsection{完整代码}
% \begin{lstlisting}
% #include <iostream>
% #include <cstdio>
% #include <algorithm>
% #define int long long
% using namespace std;

% const int INF = 1e9;
% const int MX = 5e3+5;

% int dp[4][MX]; // dp[v][c]：维生素v用c卡路里能获得的最大量
% int n, x;

% int max(int a, int b){
%     return a>b ? a : b;
% }

% signed main() {
%     // freopen("in.txt", "r", stdin);
%     scanf("%lld%lld", &n, &x);
    
%     for(int v=1; v<=3; v++) {
%         fill(dp[v], dp[v]+x+1, -INF);
%         dp[v][0] = 0; 
%     }
    
%     for(int i=1; i<=n; i++) {
%         int v, a, c;
%         scanf("%lld%lld%lld", &v, &a, &c);
%         for(int j=x; j>=c; j--) {
%             if(dp[v][j-c] != -INF) {
%                 dp[v][j] = max(dp[v][j], dp[v][j-c]+a);
%             }
%         }
%     }
    
%     for(int v=1; v<=3; v++) {
%         for(int j=1; j<=x; j++) {
%             dp[v][j] = max(dp[v][j], dp[v][j-1]);
%         }
%     }
    
%     int low = 0, high = min({dp[1][x], dp[2][x], dp[3][x]}), ans = 0;
%     if(high < 0) return printf("0"), 0;
    
%     while(low <= high) {
%         int mid = (low+high)>>1;
%         int sum = 0, ok = 1;
        
%         for(int v=1; v<=3; v++) {
%             int* d = dp[v];
%             int cost = lower_bound(d, d+x+1, mid) - d;
%             if((sum += cost) > x || cost > x) {
%                 ok = 0; break;
%             }
%         }
        
%         ok ? (ans=mid, low=mid+1) : (high=mid-1);
%     }
%     printf("%lld", ans);
% }
% \end{lstlisting}

\subsection{其实 Allen 是幕后大 BOSS}
\subsubsection{形式化定义}
给定一个长度为$n$（$n$为奇数）的员工序列，其中第$i$个员工的工资范围为$[l_i, r_i]$。初始总工资满足$\sum_{i=1}^n l_i \leq s$。我们需要为每个员工分配工资$w_i \in [l_i, r_i]$，满足：

\begin{enumerate}
    \item 总工资约束：$\sum_{i=1}^n w_i \leq s$
    \item 最大化工资序列的中位数
\end{enumerate}

设$k = \frac{n+1}{2}$，优化目标为：

$$
\max \left\{ \text{median}(w_1, w_2, \dots, w_n) \right\}
$$
约束条件：
$$
\begin{cases}
    l_i \leq w_i \leq r_i & \forall i \in \{1,2,\dots,n\} \\
    \sum_{i=1}^n w_i \leq s
\end{cases}
$$

\subsubsection{解题思路}
先发出基础工资，把每位员工工资初始化为$l[i]$，再看此时高于设定目标中位数$mid$的员工数量，若达到要求就返回，否则调整后续员工工资；为了调整中位数但尽可能小的增加总工资，应该优先调整\textcolor{blue}{最接近中位数的元素}，所以综合考虑前面的论述，我们必须要对员工的最低工资$l[i]$进行\textcolor{blue}{从大到小}排序。

当目标中位数$mid$可以达成时，则挑战更高的中位数标准，若不行，则尝试更低的中位数标准。需要注意的是，\textbf{每次尝试之间需要把总工资回调到原来的状态}，这里我们并没有对每个人的具体工资做改动，而是只改动总工资，在效果上相当于改动了每个人的具体工资。

\subsubsection{复杂度分析}
\begin{itemize}
    \item 预处理：$O(n)$计算初始总工资
    \item 排序：$O(n \log n)$
    \item 二分查找：$O(\log (\max r_i))$
    \item 每次检查：$O(n)$
    \item 总复杂度：$O(n \log n + n \log (\max r_i))$
\end{itemize}

% \subsubsection{完整代码}
% \begin{lstlisting}
% #include <bits/stdc++.h>
% #define ll long long
% using namespace std;

% const int N = 2e5 + 10;
% int t;
% int n;
% ll s;
% int l[N], r[N];
% ll ans;
% ll minl=2e14, maxr=0;


% ll min(ll a, ll b){
%     return a < b ? a : b;
% }
% ll max(ll a, ll b){
%     return a > b ? a : b;
% }
% bool cmp(int a, int b) {
%     return l[a] > l[b];
% }
% void solve() {
%     ll sum_li = 0;
%     for (int i = 1; i <= n; ++i) {
%         sum_li += l[i];
%     }
    
%     vector<int> idx(n);
%     for (int i = 0; i < n; ++i) {
%         idx[i] = i + 1;
%     }
%     sort(idx.begin(), idx.end(), cmp);
    
%     int k = (n + 1) / 2;
%     ll left = minl, right = maxr;
%     ans = 0;
    
%     while (left <= right) {
%         ll mid = left + (right - left) / 2;
%         ll current_sum = sum_li;
%         int cnt = 0;
%         int i = 0;
        
%         // 高工资统计
%         for (; i < idx.size(); ++i) {
%             int id = idx[i];
%             if (l[id] > mid) {
%                 cnt++;
%                 if (cnt == k) break;
%             } else {
%                 break;
%             }
%         }
        
%         // 高工资数量足够多，直接尝试下一种可能
%         if (cnt >= k) {
%             ans = mid;
%             left = mid + 1;
%             continue;
%         }
        
%         // 调整工资
%         for (; i < idx.size(); ++i) {
%             int id = idx[i];
%             if (r[id] >= mid) {
%                 current_sum += (mid - l[id]);
%                 cnt++;
%                 if (cnt == k) break;
%             }
%         }
        
%         if (cnt >= k && current_sum <= s) {
%             ans = mid;
%             left = mid + 1;
%         } else {
%             right = mid - 1;
%         }
%     }
% }

% int main() {
%    // freopen("in.txt", "r", stdin);
%     scanf("%d", &t);
%     while (t--) {
%         ans = 0;
%         scanf("%d%lld", &n, &s);
%         for (int i = 1; i <= n; ++i) {
%             scanf("%d%d", &l[i], &r[i]);
%             minl = min(minl, l[i]);
%             maxr = max(maxr, r[i]);
%         }
%         solve();
%         printf("%lld\n", ans);
%     }
%     return 0;
% }
% \end{lstlisting}

\subsection{Sherway与幂}
\subsubsection{形式化定义}
求1到$N$中共有多少个数可以表示成$M^K, K>1, N\le 1e18$.
\subsubsection{解题思路}
我们会先有一个朴素的思想，也就是我对$n$进行开根，把开根的结果取整后加起来：
$$
ans=\sum_{i=2}^{\infty}\lfloor(n)^{\frac{1}{i}}\rfloor
$$
但是这样我们发现，会对个别项重复计算，例如$2^6=4^3=8^2=64$被算了3次，我们应该对其进行去重。这样我们就应该联想到容斥原理来确保不会重复计算，也不会额外扣除。

设$A_i$表示所有$p[i]$次方数的集合($p[i]$表示第$i$个质数)，即：
\[
    A_i=\{x^{p[i]}\le n|x\in\mathbb{N}\}
\]
我们需要计算的就是$|\bigcup_{i=1}^{\infty}A_i|$，即所有次方数的并集。
所以有：
\begin{equation*}
    |\bigcup_{i=1}^{\infty}A_i|=\sum_{i=1}^{\infty}|A_i|
    -\sum_{1\le i<j}|A_i\cup A_j|
    +\sum_{1\le i<j<k}|A_i\cup A_j\cup A_k|\dots
\end{equation*}
为了计算上式，所以我们首先需要确定指数取值的上界，只需要找到$imax$使得$2^{p[imax]}\le n < 2^{p[imax+1]}$即可。接着在中间的计算过程中，对于由\textcolor{blue}{奇数}个质数相乘得到的指数幂集合应该\textcolor{blue}{加上}，而对于\textcolor{blue}{偶数}质数相乘得到的指数幂集合应该\textcolor{blue}{减去}。这也都是我们需要统计的。

需要额外注意的是，这里我们可以利用已求出的$imax$，那么从1到$2^{imax}$中的每一个$j$，它的二进制表示可以算作是对于素数数组的一个选取规则。例如$0b101$可以表示选择了素数2和5.这样只需要记录其中1的数量，就知道其选取的素数的奇偶数了。

\subsubsection{复杂度分析}
\begin{itemize}
    \item 外层循环：$2^{imax}-1$次
    \item 内层循环$imax$次
    \item 总复杂度$O(imax \cdot (2^{imax}-1))$(其中$imax$约为2到$\log n$中质数的数量，在本题中不超过20)
\end{itemize}


% \subsubsection{完整代码}
% \begin{lstlisting}
% #include <bits/stdc++.h>
% using namespace std;

% #define ll long long
% ll n, ans;

% long double ONE = 1.0;

% // 由于n最大是1e18，所以不会超过1 << 67
% int ps[] = {2, 3, 5, 7, 11, 13, 17, 19, 23, 29, 31, 37, 41, 43, 47, 53, 59, 61, 67};

% void solve() {
%     int max_exp = 0;
%     while(1ll << ps[max_exp + 1] < n) max_exp++;

%     for (int mk = 1; mk < (1 << max_exp); mk++) {
%         ll exp = 1;
%         bool flag = false;
        
%         for (int j = 0; j < max_exp; j++) {
%             if (mk & (1 << j)) {
%                 exp *= ps[j];
%                 flag = !flag;
%             }
%         }
%         // 特别注意此处对于浮点数精度的调整。
%         long double one = 1;
%         ll cnt = (ll)(powl(n, one / exp)+1e-10);
        
%         if (flag) {
%             ans += cnt;
%         } else {
%             ans -= cnt;
%         }
%     }
% }

% int main() {
%     scanf("%lld", &n);
%     solve();
%     printf("%lld\n", ans + 1);// 考虑到1一定满足题意，而前面没有统计。
%     return 0;
% }
% \end{lstlisting}

\end{document}